\cdpchapter{Resumen}

\begingroup

Este Trabajo Fin de Grado desarrolla un procedimiento computacional para analizar y predecir la volatilidad de datos financieros intradía. El trabajo combina una parte metodológica, una aplicación empírica sobre Bitcoin y un prototipo web que permite utilizar una parte de los modelos obtenidos.

El objetivo principal es construir un flujo de análisis reproducible, desde la obtención y preparación de los datos hasta la comparación de modelos predictivos. Para ello, primero se utiliza un caso sintético conocido, que sirve como control del procedimiento en un entorno donde la estructura de los datos está definida de antemano. Después, el mismo enfoque se aplica a datos reales de Bitcoin con frecuencia de cinco minutos, a partir de los cuales se construye una medida de volatilidad futura a corto plazo.

El trabajo no busca predecir la dirección del precio ni generar señales de compra o venta. La variable estudiada representa intensidad de movimiento, no rentabilidad esperada. Por ello, las salidas del sistema deben interpretarse como estimaciones de volatilidad, y no como recomendaciones de inversión.

La metodología desarrollada incluye preparación y validación de datos, construcción de variables, caracterización de la serie, comparación con modelos de referencia y evaluación predictiva fuera de muestra. Además, se incorpora una comprobación sintética para distinguir entre el comportamiento del procedimiento en un sistema controlado y su aplicación a una serie financiera real, donde existen ruido, persistencia y cambios de comportamiento.

Como cierre práctico, el proyecto integra los modelos y artefactos generados en un MVP web desarrollado con Streamlit. Esta aplicación permite cargar datos recientes o de ejemplo, validar su calidad, calcular las variables necesarias y mostrar predicciones de volatilidad junto con información de diagnóstico. El MVP no sustituye al análisis completo, sino que actúa como demostración funcional de una parte del trabajo realizado.

La aportación principal del TFG consiste en conectar metodología, análisis empírico e implementación software dentro de un mismo sistema reproducible. El resultado es una herramienta que permite estudiar la volatilidad de Bitcoin con un enfoque prudente, comparando distintos modelos y separando claramente lo que puede afirmarse a partir de los datos de aquello que requiere cautela interpretativa.

\endgroup